\paragraph*{04.03.01.01 Das Projekt mit der Mission und der Vision der Organisation in Einklang bringen}

%Task
\par Die Vision des \gls{adac} ist die weitere Automatisierung und Digitalisierung von Geschäftsprozessen, um die Servicequalität für Mitglieder zu steigern. Das Ziel des Programms \gls{r2h} ist es, die bestehende IT‑Landschaft durch die Einführung von \gls{s4} und weiterer Komponenten wie \gls{mbc}, \gls{aif} und \gls{drc} zu modernisieren, um die Effizienz und Flexibilität der IT‑Systeme zu erhöhen. In meiner Rolle als \gls{wsl} war es meine Aufgabe sicherzustellen, dass die strategische Ausrichtung des Programms mit der übergeordneten Mission und Vision des \gls{adac} übereinstimmt.

%Action
\par Um dies zu erreichen, wurden in der \gls{rup} fachlich vielfältige Soll‑Prozesse im \gls{gfa} definiert, die die zukünftige Arbeitsweise der Organisation widerspiegeln. Jedoch wurde die technische Betrachtungsweise in der \gls{gls-rup} anfangs vernachlässigt, was zu einem Mangel an Klarheit über die technische Umsetzung und die damit verbundenen Herausforderungen führte. In meiner Rolle als \gls{wsl} habe ich aktiv daran gearbeitet, diese Lücke zu schließen, indem ich die technischen Anforderungen und Herausforderungen in die strategische Planung integriert habe, die durch die fehlende technische Perspektive und die Projektprämisse \gls{nir} für Vorsysteme entstanden ist.

%Result
\par Um diesen Zielkonflikt zu adressieren, habe ich die technische Perspektive in die strategische Planung eingebunden, indem ich die technischen Anforderungen und Herausforderungen in die Definition der Projektprämissen aufgenommen habe. Dies führte dazu, dass die Projektprämisse \gls{nir} für nicht am \gls{eol} befindliche Vorsysteme angepasst werden musste, wenn die technologische Umsetzung im neuen Framework \gls{aif} erfolgen soll. Der Schwenk auf eine \gls{cc} zuzüglich der notwendigen \gls{s4}-Enhancements, welche den technologischen und fachlichen Veränderungen für die bestehenden \gls{ricefw}-Objekte Rechnung tragen, erfolgte, um das Termin- und Qualitätsziel der ersten Migrationswelle zu erreichen. Die daraus resultierende \gls{td} war die Folge.

%Reflect
\par Es ist entscheidend, dass die technische Umsetzung zwingend in die fachlichen Konzeptionsworkshops einbezogen wird, um die strategische Ausrichtung des Projekts mit der technologischen Realität in Einklang zu bringen. Die frühzeitige Einbindung der technischen Perspektive ermöglicht es, potenzielle Herausforderungen und Risiken rechtzeitig zu identifizieren und proaktiv anzugehen, um sicherzustellen, dass die Projektziele erreicht werden können. Dies wird für die Welle 1.3 ff. berücksichtigt.
